\documentclass[a4paper, 12pt]{report}

% --- Configuração de Idioma e Tipografia ---
\usepackage[T1]{fontenc}      % Codificação de fonte moderna (essencial para hifenação e PDF)
\usepackage[utf8]{inputenc}   % Permite escrever acentos diretamente
\usepackage[portuguese]{babel}% Suporte ao idioma português (hifenação, etc.)
\usepackage{csquotes}
\usepackage{lmodern}          % Fonte Latin Modern (melhor que a padrão)
\usepackage{indentfirst}      % Indenta o primeiro parágrafo de cada seção
\usepackage{url}

% --- Pacotes de Layout e Estilo ---
\usepackage[margin=20mm]{geometry} % Define as margens do documento
\usepackage{graphicx}             % Para incluir imagens (o padrão moderno)
\usepackage{setspace}             % Para controle de espaçamento (\singlespacing, etc)
\usepackage{fancybox}             % Para caixas estilosas como \ovalbox
\usepackage[table]{xcolor}        % Para usar cores, especialmente em tabelas
\usepackage{caption}              % Para customizar legendas de figuras e tabelas
\usepackage{titlesec}             % Para customizar títulos de capítulos e seções
\usepackage{float}                % Para controlar a posição de figuras e tabelas 

% --- Pacotes para Conteúdo Específico ---
\usepackage{listings}             % Para blocos de código simples
\lstset{
    inputencoding=utf8,
    extendedchars=true,
    literate=
        {á}{{\'a}}1 {ã}{{\~a}}1 {â}{{\^a}}1 {à}{{\`a}}1
        {é}{{\'e}}1 {ê}{{\^e}}1
        {í}{{\'i}}1
        {ó}{{\'o}}1 {õ}{{\~o}}1 {ô}{{\^o}}1
        {ú}{{\'u}}1
        {ç}{{\c{c}}}1
        {Á}{{\'A}}1 {Ã}{{\~A}}1 {Â}{{\^A}}1 {À}{{\`A}}1
        {É}{{\'E}}1 {Ê}{{\^E}}1
        {Í}{{\'I}}1
        {Ó}{{\'O}}1 {Õ}{{\~O}}1 {Ô}{{\^O}}1
        {Ú}{{\'U}}1
        {Ç}{{\c{C}}}1
}
\usepackage{longtable}            % Para tabelas que podem quebrar entre páginas
\usepackage{tikz}                 % Para criar diagramas e gráficos vetoriais
\usetikzlibrary{arrows.meta, positioning, shapes.geometric, fit}

% --- Configuração da Bibliografia (BibLaTeX) ---
\usepackage[backend=biber, style=numeric, sorting=none]{biblatex}
\addbibresource{referencias.bib}

% --- Pacotes de Links e URLs ---
\usepackage{xurl}                 % Permite quebras de linha em qualquer lugar de uma URL
\usepackage{hyperref}             % Torna referências e URLs clicáveis
\hypersetup{hidelinks}            % Remove as caixas coloridas dos links

\newcommand{\chapterauthor}[1]{
  {\large\textbf{Autor(es):} \textit{#1}\par\vspace{20pt}}
}

\newcommand{\chapterrev}[1]{
  {\large\textbf{Revisão:} \textit{#1}\par\vspace{20pt}}
}

\titleformat{\chapter}[display]
  {\normalfont\Huge\bfseries}
  {}
  {0pt}
  {}

\titlespacing*{\chapter}
  {0pt}
  {-20pt}
  {15pt}



\definecolor{codegreen}{rgb}{0,0.6,0}
\definecolor{codegray}{rgb}{0.5,0.5,0.5}
\definecolor{codepurple}{rgb}{0.58,0,0.82}
\definecolor{backcolour}{rgb}{0.95,0.95,0.92}

\lstdefinestyle{mystyle}{
    backgroundcolor=\color{backcolour},   
    commentstyle=\color{codegreen},
    keywordstyle=\color{magenta},
    numberstyle=\tiny\color{codegray},
    stringstyle=\color{codepurple},
    basicstyle=\ttfamily\footnotesize,
    breakatwhitespace=false,         
    breaklines=true,                 
    captionpos=b,                    
    keepspaces=true,                 
    numbers=left,                    
    numbersep=5pt,                  
    showspaces=false,                
    showstringspaces=false,
    showtabs=false,                  
    tabsize=2
}


\lstdefinelanguage{arm}{
    morekeywords={sub, stmfd, blx, ldmfd, mov, add, cmp, b, bl, ldr, mcr, mrs, bic, orr, msr, bx},
    sensitive=false,
    morecomment=[l]{//}, % recognize // as comments
    morecomment=[l]{@},  % recognize @ as comments
    morestring=[b]",
    basicstyle=\ttfamily\small,
    keywordstyle=\color{blue}\bfseries,
    commentstyle=\color{green!60!black},
}

\lstset{style=mystyle}



\title{Kernel: o Núcleo do Sistema Operacional}
\author{Comunidade UFSKernel}
\date{\today}

\begin{document}

\input{capa.tex}

\tableofcontents
\newpage


\chapter{Infraestrutura Base}

\chapterauthor{Ti Ribeiro Silvério}

%-------------------------------------------------------------------------------
\section{Descrição}
%-------------------------------------------------------------------------------
% Descreva em alto nível o que este módulo faz e qual problema ele resolve
% dentro do sistema operacional.
%

Esse é o kernel mínimo para possibilitar o desenvolvimento dos módulos subsequentes, um executável para o \verb|QEMU ARM| que coloca o processador para executar código \verb|C| compilado.

%-------------------------------------------------------------------------------
\section{Decisões de Arquitetura e Design}
%-------------------------------------------------------------------------------
% Justificar o porque de ter feito as coisas de uma certa maneira.
%
% Tópicos a cobrir:
% - Quais alternativas foram consideradas?
% - Por que a abordagem escolhida foi a melhor para este projeto?
% - Quais são as limitações ou trade-offs da abordagem escolhida?
%

Para o desenvolvimento deste projeto, foi adotada uma filosofia de simplicidade máxima. Nele, foi feito o mínimo possível para que o processador ARM virtual executasse código \verb|C| compilado. Essa escolha foi feita para tornar o entendimento dos detalhes mais fácil, e possibilitar que funcionalidades extras sejam adicionadas à medida que forem necessárias.

Um exemplo disso foi a preparação da pilha apontada pelo registrador \verb|sp|. Para solucionar esse problema, foi alocado um espaço de 32KB no binário final dentro do \textit{linker script}. Porém, por mais que tenha sido definido um rótulo que armazena a posição de memória do fim da pilha, não foi implementada uma maneira de tratar o estouro. Dessa forma, essas soluções serão implementadas em projetos futuros à medida que problemas aparecerão.

%-------------------------------------------------------------------------------
\section{Detalhes da Implementação}
%-------------------------------------------------------------------------------
% Descreva como a solução foi implementada, conectando o código
% à teoria.
%
% Tópicos a cobrir:
% - Principais estruturas de dados (structs) e suas funções.
% - Algoritmos chave utilizados.
% - Fluxo de execução das funções mais importantes.
% - Inclua pequenos trechos de código com o pacote 'listings'.
%

Esse projeto consiste em 4 partes:

\begin{enumerate}
    \item O código de inicialização em Assembly ARM;
    \item O \textit{linker script} que liga o código de inicialização ao código em \verb|C| compilado e prepara a pilha;
    \item A função principal em \verb|C| para a qual o processador é desviado no código de inicialização.
\end{enumerate}

\subsection{Código de inicialização}

O código em Assembly é o mais simples de todos. Primeiro, ele coloca o processador no modo supervisor e desliga interrupções, depois configura a pilha e desvia a execução para o código em \verb|C|.

\begin{lstlisting}[caption=Configuração do modo do processador\, da pilha e chamada a kmain().]
_start:
    mrs r0, cpsr
    bic r0, r0, #0x1f
    orr r0, r0, #0xd3
    msr cpsr, r0

    ldr sp, =stack_top
    bl kmain

hang:
    b hang
\end{lstlisting}

Nesse código, o registrador \verb|r0| é usado como registrador auxiliar para executar as mudanças. Assim, é escrito $11010011_2$ no útimo byte do registrador \textit{cpsr}, que armazena o estado atual do modo de operação do processador. Os dois bits mais a esquerda fazem o processador ignorar interrupções, enquanto os cinco mais a direita ligam o modo supervisor \cite{arm:armv7ar}.

É importante desativar as interrupções pois ainda não definimos como tratá-las. Caso acontecesse qualquer interrupção durante a execução de um programa, sua execução seria desviada para uma posição de memória desconhecida e isso travaria o sistema.

Também é essencial ligar o modo supervisor, pois o kernel deve ser capaz de ter controle total sobre o computador e sua memória, diferente dos demais processos que rodarão no modo usuário.

Após isso, é configurada a pilha e desviado o contador de programa para o binário da função \verb|kmain()|. Caso haja um retorno, o processador entra em um laço infinito para impedir a execução de lixo de memória.

A inicialização do \textit{stack pointer} é de extrema importância, pois chamadas de funções e passagem de parâmetros são compiladas para saltos e empilhamento de valores. Sem a existência da pilha, não seria possível chamar funções em \verb|C|.

\subsection{\textit{Linker script}}

Geralmente, quando usamos o \textit{linker} para terminar de gerar um executável, nem nos preocupamos com \textit{como} os binários serão ligados. No caso do kernel, esse não é o caso, e precisamos organizar as seções dos binários a serem ligados de tal forma que sejam executadas corretamente durante o \textit{boot}. Para resolver esse problema, temos que definir essas especificidades em uma receita para o \textit{linker}.

Esse \textit{script}, assim como o código de inicialização, é bem simples.

\begin{lstlisting}[language=arm, caption=Script de geração do binário final.]
ENTRY(_start)

SECTIONS {
    . = 0x40000000;

    .text : {
        KEEP(*(.text.startup))
        *(.text)
        *(.rodata)
    }

    .data : {
        *(.data)
    }

    .bss : {
        *(.bss)
    }

    stack_bottom = .;
    . = . + 0x8000;
    stack_top = .;

    end = .;
}
\end{lstlisting}

A primeira linha é responsável por apontar que o início lógico do executável deve ser \verb|_start|, o rótulo definido no código Assembly.

O bloco da linha 3 é a definição da ordem dos binários a serem ligados no binário final. Nesse bloco, definimos o início da DRAM (endereço \verb|0x40000000| \cite{qemu:virtmachine}) e o que vai dentro de cada seção do executável: \verb|.text|, \verb|.data| e \verb|.bss|.

A seção \verb|.text| é onde deve entrar o código fonte do programa, pois é uma seção com permissões de leitura e execução. Nela, colocamos o código de iniciação da subseção anterior, as seções \verb|.text| de cada binário a ser ligado e as seções \verb|.rodata| de cada um deles.

Essa ultima seção (\verb|.rodata|) não armazena código, mas sim valores constantes, como cadeias de caracteres literais como em \texttt{print("Hello World")}. A constante \texttt{"Hello World"} não é uma variável, mas ainda precisa ser armazenada em algum lugar, nesse caso, \verb|.rodata|. Como \verb|.text| é uma seção de apenas leitura, é um bom lugar para armazenar esse tipo de dado.

Similarmente, as seções \verb|.data| e \verb|.bss|, que contém as variáveis globais, são preenchidas com a junção das seções de mesmo nome dos arquivos objeto.

Por fim, vem a definição das posições de inicio e fim da pilha e a alocação de seu espaço. O fim\footnote{Observação: a pilha ``cresce para baixo'', com o topo começando em um valor alto e decrescendo até um valor mais baixo. Decidi, por clareza, chamar de ``fim'' o ponto mais baixo, o limite de decrescimento, e de ``início'' o ponto mais alto, onde ela começa vazia.} da pilha é definido como o endereço logo após a seção \verb|.bss|, são pulados 32KB, e o início da pilha é endereço após esse salto. Dessa forma, reservamos um espaço para que a pilha possa funcionar normalmente.

\subsection{Código em \texttt{C}}

A função \texttt{kmain}, em \texttt{C}, é exatamente para onde o código em Assembly salta. É importante lembrar que tudo que acontece aqui é em modo supervisor. Desse ponto em diante, começa o trabalho dos grupos. Módulos são inicializados, testados, processos são iniciados e etc.

\begin{lstlisting}[language=C, caption=Função principal.]
#include <serial.h>

void kmain(void) {
    // ...
    // Inicializações e testes dos módulos
    // ...

    first_process(prog_init_idle);

    return;
}
\end{lstlisting}

No fim da função, o kernel executa seu primeiro processo (mais sobre processos a frente), em modo usuário. Esse processo inicial é responsável por iniciar a Shell (outro processo) e existir sem fazer nada, ou seja, um \textit{idle process}. A existência do \textit{idle process} é importante para que sempre haja um processo para o escalonador escolher, possibilitando assim o tratamento de interrupções, que são desabilitadas no modo supervisor.

\chapter{Tratador de Interrupções}

\chapterauthor{Renan Matulianes Arnaldo, Renan Suana Grothe Garcia, João Vitor Grisoto Kussaba}

%-------------------------------------------------------------------------------
\section{Descrição}
%-------------------------------------------------------------------------------
% Descreva em alto nível o que este módulo faz e qual problema ele resolve
% dentro do sistema operacional.
%
O módulo de tratador de interrupções de um Sistema Operacional (SO) é responsável por receber as interrupções geradas pelo \emph{hardware} ou pelo \emph{software} e executar a respectiva rotina de tratamento para o tipo de interrupção recebida.

As interrupções são sinais gerados pelo \emph{hardware} ou \emph{software} que indicam que algum evento precisa de atenção completa do processador para serem tratados corretamente. Essas interrupções podem ocorrer por meio de chamadas de sistema, instruções específicas, erros em tempo de execução, entre outras rotinas. Desse modo, o tratador de interupções é um módulo essencial do \emph{kernel} do SO, porque ele define as rotinas de tratamento para realizar cada uma das respostas apropriadas para os sinais de interrupção recebidos.


%-------------------------------------------------------------------------------
\section{Decisões de Arquitetura e Design}
%-------------------------------------------------------------------------------
% Justificar o porque de ter feito as coisas de uma certa maneira.
%
% Tópicos a cobrir:
% - Quais alternativas foram consideradas?
% - Por que a abordagem escolhida foi a melhor para este projeto?
% - Quais são as limitações ou trade-offs da abordagem escolhida?
%
Para a implementação do tratador de interrupções, foi utilizada uma abordagem similar ao SO Minix, que é a utilização de uma tabela de vetores de bits que endereça fisicamente na memória onde está localizado a rotina espefífica de tratamento de interrupção, onde cada entrada dessa tabela possui um tipo de interrupção endereçada pelo \emph{offset} de bits enviados.

Essa abordagem foi escolhida para que o acesso à rotina de tratamento de interrupções possa ser feito de forma constante, aumentando a eficiência em que o tratamento de interrupções ocorre e reduzindo o período em que o precessador mantêm-se ocupado com a rotina de tratamento.
%-------------------------------------------------------------------------------
\section{Detalhes da Implementação}
%-------------------------------------------------------------------------------
% Descreva como a solução foi implementada, conectando o código
% à teoria.
%
% Tópicos a cobrir:
% - Principais estruturas de dados (structs) e suas funções.
% - Algoritmos chave utilizados.
% - Fluxo de execução das funções mais importantes.
% - Inclua pequenos trechos de código com o pacote 'listings'.
%
Seguindo a linha de pensamento abordada, foi utilizado um vetor de registro de funções para a implementação dos handlers de interrupções, que utiliza de ponteiros de função para acessar as rotinas de tratamento de interrupção. Além disso, é utilizada uma \emph{struct} denominada \emph{Interrupts Sevice Routine} (isr) responsável por realizar o resgistro das funções que farão o tratamento de interrupções pelos dispositivos.

Para isso, são utilizadas algumas funções tem tem por objetivo a manipulação dos bits presentes nos registradores GIC (\emph{Generic Interrupt Controller}) da arquitetura ARM (arquitetura selecionada para a execução do projeto), que são registradores divididos por blocos lógicos responsáveis por gerenciar, priorizar e rotear interrupções. Para a etapa realizada neste projeto, os blocos mais utilizados para a configuração de interrupções no núcleo do sistema operacional são o bloco distribuidor \emph{Generic Interrupt Controller Distributor} (GICD) e o bloco de interface da CPU \emph{Generic Interrupt Controller CPU Interface} (GICC).

Esses blocos são manipulados pelas funções \emph{gic\textunderscore{}enable\textunderscore{}interrupt()}, responsável por habilitar e registrar interrupções e modificar os registradores GIC\textunderscore{}ISENABLER para que permitam o envio da interrupção à CPU;
\begin{lstlisting}
void gic_enable_interrupt(uint32_t irq) {
  if (irq >= MAX_INTERRUPTS)
    return;

  // cada registrador tem 32 bits
  uint32_t reg_index = irq / 32;  // qual registrador
  uint32_t bit_offset = irq % 32; // qual bit do registrador

  GICD_ISENABLER[reg_index] |= (1 << bit_offset);
}
\end{lstlisting}
e \emph{gic\textunderscore{}disable\textunderscore{}interrupt()}, análogo ao \emph{gic\textunderscore{}enable\textunderscore{}interrupt()}, só que desativa a interrupção selecionada e modifica os registradores GICD\textunderscore{}ICENABLER para sinalizar que a interrupção está desativada;
\begin{lstlisting}
void gic_disable_interrupt(uint32_t irq) {
  if (irq >= MAX_INTERRUPTS)
    return;

  uint32_t reg_index = irq / 32;
  uint32_t bit_offset = irq % 32;

  // No ICENABLER escrever 1 desliga a interrupção
  GICD_ICENABLER[reg_index] = (1 << bit_offset);
}
\end{lstlisting}

Desse modo, o fluxo principal de execução funciona da seguinte forma: Um gatilho de \emph{hardware} altera os bits de informação que são lidos pelo tratador como um tipo volatile uint32\textunderscore{}t (sendo o tipo uint32\textunderscore{}t um \emph{typedef} do tipo unsigned int da linguagem C), que indica ao compilador que o valor de bits pode ser mudado inesperadamente fora do controle ou da visibilidade da execução do tratador, e são levados aos registradores GIC;
\begin{lstlisting}
// Registradores
#define GICD_CTLR (*(volatile uint32_t *)(GICD_BASE + 0x000))
#define GICC_CTLR (*(volatile uint32_t *)(GICC_BASE + 0x000))
#define GICC_PMR (*(volatile uint32_t *)(GICC_BASE + 0x004))
#define GICC_IAR (*(volatile uint32_t *)(GICC_BASE + 0x00C))
#define GICC_EOIR (*(volatile uint32_t *)(GICC_BASE + 0x010))
\end{lstlisting}
após isso, a informação é levada à tabela de vetores em linguagem de máquina, que desvia a execução do programa para ela;
\begin{lstlisting}
vector_table:			// offset |  interrupcao  |  modo de execucao 
    b _start			// 0x00   |   Reset	  |  supervisor
    b undefined_handler     	// 0x04   |   Undefined	  |  undefined
    b svc_handler           	// 0x08   |   SVC  	  |  supervisor
    b .                     	// 0x0C   |   Prefetch    |  abort
    b data_abort_handler    	// 0x10   |   Data Abort  |  abort 
    b .                     	// 0x14   |   Hyp Trap    |  hypervisor
    b irq_handler           	// 0x18   |   IRQ         |  interrupt
    b .                     	// 0x1C   |   FIQ         |  fast interrupt
\end{lstlisting}
depois, o endereço da função em linguagem de máquina \emph{irq\textunderscore{}handler} ajusta o resgistrador de retorno da arquitetora para o processo o qual estava em operação antes da interrupção ter sido gerada e empilha todos os registradores na pilha presente na memória do processador para salvar o contexto antes de ocorrer a troca para o núcleo do SO;
\begin{lstlisting}
irq_handler:
    sub lr, lr, #4
    
    // Uma pseudo-instrução que armazena os registradores de uso geral R0-R12 e o LR (Link Register)
    // e ao mesmo tempo atualiza o valor do stack pointer (sp) para o topo da pilha após o armazenamento stmfd sp!, {r®-r12, lr}
    blx irq_dispatcher_c // Salta para a função em C
    
    // Uma pseudo-instrução que restaura os registradores R0-R12 da pilha
    // e ao mesmo tempo copia o valor antigo de LR direto para o program counter (PC)
    ldmfd sp!, {rø-r12, pc}^
\end{lstlisting}
uma vez que o tratador de interrupções está em execução, o registrador GICC\textunderscore{}IAR é lido para que o tratador saiba que interrupção o chamou e se essa interrupção possui uma rotina de tratamento configurada para ela. Caso possua, o tratador faz o desvio para a rotina específica armazenada no vetor de ponteiros de funções \emph{interrupt\textunderscore{}handlers} e salva o identificador da rotina no resgistrador GICC\textunderscore{}EOIR, que indica que a interrupção já foi servida;
\begin{lstlisting}
void irq_dispatcher_c(void) {
  uint32_t irq_id = GICC_IAR & 0x3FF; // lê o identificador da interrupção

  // verifica se o indenficador não é maior que o máximo de interrupções configuradas e se existe uma rotina de tratamento armazenada no vetor interrupt_handlers
  if (irq_id < MAX_INTERRUPTS && interrupt_handlers[irq_id] != 0) {
    interrupt_handlers[irq_id]();
  }

  GICC_EOIR = irq_id; // marca que a interrupção já foi servida
}
\end{lstlisting}
por fim, o tratador realiza um retorno para o contexto do processo anterior a sua chamada para que ele possa continuar com a sua execução.
%-------------------------------------------------------------------------------
\section{Interface Pública (API)}
%-------------------------------------------------------------------------------
% Documente como outros módulos do kernel podem usar este módulo.
% Essencialmente, é a documentação do arquivo .h correspondente.
%
% Para cada função pública:
% - Nome da função.
% - Descrição do que ela faz.
% - Parâmetros (o que cada um significa).
% - Valor de retorno (o que ela retorna e em que casos).
%
Para a utilização e configuração do tratador de forma que seja possível configurar e gerenciar interrupções, a biblioteca desenvolvida possui as seguintes funções:
\begin{itemize}
    \item int register\textunderscore{}interrupt\textunderscore{}handler(uint32\textunderscore{}t irq, isr\textunderscore{}t handler)

Recebe o identificador de interrupção \emph{irq} e um ponteiro de função \emph{handler} que aponta para a rotina de tratamento da interrupção. Essa função salva o ponteiro de função no vetor \emph{interrupt\textunderscore{}handlers} no índice \emph{irq} caso o identificador recebido não ultrapasse o valor MAX\textunderscore{}INTERRUPTS. Retorna 0 caso tenha sido possível salvar o \emph{handler} senão -1.
    \item void gic\textunderscore{}enable\textunderscore{}interrupt(uint\textunderscore{}t irq)

Recebe o identificador de interrupção \emph{irq} e habilita o registrador GICD\textunderscore{}ISENABLER correspondente caso identificador não ultrapasse o valor MAX\textunderscore{}INTERRUPTS, habilitando-o para o tratamento de interrupção adequado.
    \item void gic\textunderscore{}disable\textunderscore{}interrupt(uint\textunderscore{}t irq)

Recebe o identificador de interrupção \emph{irq} e habilita o registrador GICD\textunderscore{}ICENABLER correspondente caso identificador não ultrapasse o valor MAX\textunderscore{}INTERRUPTS, desabilitando o tratamento de interrupção correspondente.
    \item void gic\textunderscore{}config\textunderscore{}interrupt(uint32\textunderscore{}t irq, int edge\textunderscore{}triggered)
Recebe o identificador de interrupção \emph{irq} e um sinal \emph{edge\textunderscore{}triggered} que indica se a interrupção é sensitiva à mudança de estado imediata (valor 1) ou somente após a estabilização do sinal recebido (valor 0) e configura o índex do identificador no registrador GICD\textunderscore{}ICFGR correspondente.
\end{itemize}
%-------------------------------------------------------------------------------
\section{Testes e Verificação}
%-------------------------------------------------------------------------------
% Mostre como o funcionamento do módulo foi testado e validado.
%
% Tópicos a cobrir:
% - Descreva a rotina de teste que foi adicionada à função kmain().
% - Inclua a saída esperada (e obtida) na console serial do QEMU.
% - Explique quais casos de borda foram testados.
%
Para o teste da funcionalidade, foi implementado o tratador para a interrupção de timer no método principal do projeto \emph{kmain()}, que tem por objetivo imprimir na tela uma mensagem de aviso que o processador sofreu uma interrupção de timer.

Para isso, na função \emph{kmain()}, é realizado a inicialização do serial do projeto por meio da função \emph{serial\textunderscore{}init()}, depois a CPU é inicializda para receber as estruturas de tratamento de interrupção por meio da função \emph{setup\textunderscore{}core\textunderscore{}for\textunderscore{}irq}, depois de inicializadas as estruturas para o tratamento de interrupções, as interrupções de CPU são ativadas com a função \emph{enable\textunderscore{}cpu\textunderscore{}interrupts()} e, por fim, é chamada a função de inicialização e configuração do timer \emph{init\textunderscore{}timer}. Após essas configurações, um laço infinito é definido que imprime uma mensagem no console do QEMU indicando que o processo está operando normalmente no processador, com um segundo laço que serve apenas para espaçar o tempo de impressão de uma mensagem para outra, até sofrer uma interrupção de timer.
\begin{lstlisting}
void kmain(void) {
  serial_init();
  serial_puts("Bem-vindo ao UFSKernel!\n");
  serial_puts("Executando em modo ARM bare-metal no QEMU.\n");

  serial_puts("Preparando CPU (VBAR e Pilha IRQ)...\n");
  setup_core_for_irq();

  serial_puts("Configurando GIC...\n");
  init_gic();

  serial_puts("Ligando interrupcoes...\n");
  enable_cpu_interrupts();

  serial_puts("Ativando Timer...\n");
  init_timer();

  while (1) {
    serial_puts("Kernel rodando na main...\n");
    // Um delay apenas para não flodar a tela
    for (volatile int i = 0; i < 50000000; i++);
  }
}
\end{lstlisting}

O inicializador do timer \emph{init\textunderscore{}timer()} lê a frequência de ciclo do hardware, divide a frequência de ciclo pelo valor SYSTEM\textunderscore{}HZ, que marca a frequência das interrupções de timer, para indicar quantos ticks ele deve aguardar antes de realizar uma interrupção, inicializa o timer com o valor de ticks calculado, registra a interrupção pela função \emph{register\textunderscore{}interrupt\textunderscore{}handler()} e habilita a interrupção pela função \emph{gic\textunderscore{}enable\textunderscore{}interruption()}.
\begin{lstlisting}
void init_timer(void) {
  uint32_t hardware_freq = read_cntfrq();
  // quantos ticks formam um quantum no sistema
  ticks_p_q = hardware_freq / SYSTEM_HZ;

  // inicializa o timer com ticks_p_q
  __asm__ volatile("mcr p15, 0, %0, c14, c2, 0" ::"r"(ticks_p_q));
  uint32_t ctl = 1;
  __asm__ volatile("mcr p15, 0, %0, c14, c2, 1" ::"r"(ctl));

  // Toda interrupção deve se registrar no despachante do GIC
  if (!register_interrupt_handler(TIMER_IRQ, timer_callback))
    serial_puts("Registrando timer no GIC...\n");
  else
    serial_puts("Erro ao registrar timer no GIC!\n");

  // ativando a interrupção apos registrar
  gic_enable_interrupt(TIMER_IRQ);
}
\end{lstlisting}
Para a rotina de testes da interrupção configurada, a interrupção de timer chama a função \emph{timer\textunderscore{}callback()}, que imprime uma pequena mensagem alertando que uma interrupção de timer ocorreu, exatamente como deveria, variando apenas caso o valor de SYSTEM\textunderscore{}HZ for configurado de uma forma diferente da já definida para o processo.
\begin{lstlisting}
// função que será chamada quando a interrupção for ativada
static void timer_callback(void) {
  serial_puts(">>> TICK DO TIMER <<<\n");

  // recarrega o timer
  __asm__ volatile("mcr p15, 0, %0, c14, c2, 0" ::"r"(ticks_p_q));

  // é aqui que o código fará a chamada para o
  // escalonador
}
\end{lstlisting}
\chapter{Entrada e Saída de teclado}

\chapterauthor{Matteo Cileneo Savan, Pedro Achilles Domingues Alves Pereira e Raul Yoshiyuki Komai}

%-------------------------------------------------------------------------------
\section{Descrição}
%-------------------------------------------------------------------------------

A entrada e saída (E/S) de dados formam a base da utilidade de qualquer sistema computacional. Sem esses subsistemas, a máquina torna-se incapaz de se comunicar com o mundo exterior. Um computador isolado é, na prática, apenas um amontoado de silício incapaz de resolver problemas reais. Neste contexto, esta seção detalha a implementação dos mecanismos de E/S, com foco na captura de dados via teclado e na saída visual em tela.

Tradicionalmente, a arquitetura dos dispositivos de E/S pode ser dividida em dois componentes principais:

\begin{description}
    \item[Componente Físico (ou Mecânico)] o dispositivo em si. É o hardware responsável por executar a operação de movimentar, exibir ou capturar os dados (as teclas de um teclado ou o painel de um monitor);
    \item[Componente Eletrônico] a controladora de dispositivo, que atua como interface entre o sistema principal e o componente físico. Na prática, ela expõe registradores para a comunicação com a CPU; a partir desses registradores, o sistema operacional torna-se capaz de enviar e receber dados, ler o estado do dispositivo, entre outras operações de controle.
\end{description}

O foco da implementação deste módulo é interagir com as controladoras de dispositivo por meio da criação de um \textit{driver} orientado a interrupções, eliminando a ineficiência da espera ocupada. Além disso, abordaremos o gerenciamento de \textit{buffers} de E/S e a implementação de interatividade com o usuário por meio de um \textit{shell} básico para a execução de comandos.

No ambiente do QEMU, a controladora simulada é a PrimeCell UART (PL011), que serve como base para a comunicação serial entre os dispositivos de E/S e o \textit{kernel}. É importante ressaltar que, neste cenário, não ocorre comunicação direta da UART com o teclado físico, tampouco com o monitor; toda essa operação é intermediada pelo próprio emulador QEMU \cite{ekasiswanto:qemu-uart}. Dessa forma, a plataforma expõe a interface necessária para a manipulação da controladora em baixo nível.


%-------------------------------------------------------------------------------
\section{Decisões de Arquitetura e Design}
%-------------------------------------------------------------------------------

O projeto visa tornar o sistema mais modular e facilitar futuras manutenções. Para refletir essa decisão estrutural, a implementação foi dividida em dois módulos: o arquivo \texttt{uart.c} concentra as funções de baixo nível da controladora UART, enquanto o arquivo \texttt{kstdio.c} fornece uma camada de abstração de mais alto nível, focada nas rotinas de E/S.

A respeito do gerenciamento temporário de dados na implementação da UART, optou-se pela utilização de um \textit{buffer} circular. Essa estrutura de dados foi escolhida por permitir o armazenamento e o consumo contínuo de fluxos de dados de maneira eficiente, operando sobre um espaço de memória estático e pré-alocado, além de apresentar uma implementação direta e de baixa complexidade.

Por fim, a arquitetura de comunicação com os dispositivos baseia-se estritamente em um sistema orientado a interrupções. Essa decisão evita o desperdício de recursos computacionais: o processador fica livre para executar outros processos ou tarefas do sistema, sendo notificado pelo hardware apenas quando ocorre, de fato, um evento (como o pressionamento de uma tecla), eliminando a ineficiência da espera ocupada.


%-------------------------------------------------------------------------------
\section{Detalhes da Implementação}
%-------------------------------------------------------------------------------

\subsection{Arquivos e Organização}

A base de código foi estruturada nos seguintes arquivos:

\begin{description}
    \item[\texttt{kstdio.c}] camada de abstração de alto nível sobre as funcionalidades da UART. Fornece um conjunto de rotinas de entrada e saída baseadas nas funções primitivas de \texttt{uart.c}, sendo amplamente utilizada no interpretador de comandos do sistema;
    \item[\texttt{main.c}] ponto de entrada do \textit{kernel}, coordena a inicialização geral do sistema;
    \item[\texttt{shell.c}] \textit{shell} interativo básico, capaz de ler, interpretar e executar comandos do usuário a partir do terminal do QEMU;
    \item[\texttt{uart.c}] encapsula toda a lógica de implementação do \textit{driver} da controladora UART emulada.
\end{description}

\subsection{Fluxo Simplificado de Execução}

Podemos dividir a execução das operações de E/S em dois fluxos principais: o de entrada (leitura) e o de saída (escrita). A seguir, detalhamos o caminho percorrido pelos dados ao atravessarem as diferentes camadas de abstração do sistema.

\subsubsection{Fluxo de Entrada}

Neste fluxo, ilustrado na Figura \ref{fig:fluxo-in}, o dado nasce no mundo físico e viaja até ser processado pelo sistema operacional. O usuário digita caracteres no teclado, que são capturados pelo terminal do QEMU e encaminhados para a UART emulada. Ao receber os dados, a UART aciona uma interrupção de hardware. O \textit{driver} intercepta essa interrupção, extrai os caracteres e os armazena de forma segura em um \textit{buffer} de entrada, disponibilizando-os para o \textit{shell} por meio de funções de leitura, como \texttt{uart\_getc()}.

\begin{figure}[htbp]
    \centering
    \resizebox{\textwidth}{!}{%
    \begin{tikzpicture}[
        >=stealth,
        box/.style={
            rectangle, 
            draw=black!70, 
            thick, 
            rounded corners=3pt,
            minimum height=1.1cm,
            text width=2.4cm,
            align=center, 
            fill=white, 
            font=\scriptsize\sffamily
        },
        arrow/.style={
            ->, 
            thick, 
            draw=black!70
        }
    ]

    \filldraw[fill=gray!5, draw=black!30, dashed, rounded corners] (2.4, -1.2) rectangle (9.6, 1.4);
    \node[font=\itshape\scriptsize, text=black!60, anchor=south] at (6.0, 0.9) {Ambiente Emulado (QEMU)};

    \filldraw[fill=blue!5, draw=black!30, dashed, rounded corners] (11.4, -1.2) rectangle (18.6, 1.4);
    \node[font=\itshape\scriptsize, text=black!60, anchor=south] at (15.0, 0.9) {Espaço do Kernel};

    \node[box] (in_user)   at (0, 0)    {Teclado Físico\\(Usuário)};
    \node[box] (in_term)   at (4.0, 0)  {Terminal Serial\\(Captura)};
    \node[box] (in_uart)   at (8.0, 0)  {UART PL011\\(Hardware)};
    \node[box] (in_driver) at (13.0, 0) {Driver UART\\(Tratamento)};
    \node[box] (in_shell)  at (17.0, 0) {Shell\\(Processa)};

    \draw[arrow] (in_user) -- (in_term);
    \draw[arrow] (in_term) -- (in_uart);
    \draw[arrow] (in_uart) -- node[above=2pt, font=\scriptsize, text=black!80] {Interrupção} (in_driver);
    \draw[arrow] (in_driver) -- node[above=2pt, font=\scriptsize, text=black!80] {\texttt{getc()}} (in_shell);

    \end{tikzpicture}%
    }
    \caption{Pipeline do fluxo de entrada.}
    \label{fig:fluxo-in}
\end{figure}

\subsubsection{Fluxo de Saída}

No fluxo de saída, ilustrado na Figura \ref{fig:fluxo-out}, o caminho se inverte. O \textit{shell} utiliza funções da camada de abstração (como \texttt{kprintf()}, \texttt{kputs()} e \texttt{kputc()}) para enviar respostas ao \textit{driver} da UART. O \textit{driver}, por sua vez, escreve esses caracteres nos registradores da controladora emulada, que os transmite de volta ao QEMU para serem finalmente renderizados no terminal do usuário.

\begin{figure}[htbp]
    \centering
    \resizebox{\textwidth}{!}{%
    \begin{tikzpicture}[
        >=stealth,
        box/.style={
            rectangle, 
            draw=black!70, 
            thick, 
            rounded corners=3pt,
            minimum height=1.1cm,
            text width=2.4cm,
            align=center, 
            fill=white, 
            font=\scriptsize\sffamily
        },
        arrow/.style={
            ->, 
            thick, 
            draw=black!70
        }
    ]

    \filldraw[fill=blue!5, draw=black!30, dashed, rounded corners] (-1.6, -1.2) rectangle (5.6, 1.4);
    \node[font=\itshape\scriptsize, text=black!60, anchor=south] at (2.0, 0.9) {Espaço do Kernel};

    \filldraw[fill=gray!5, draw=black!30, dashed, rounded corners] (6.4, -1.2) rectangle (13.6, 1.4);
    \node[font=\itshape\scriptsize, text=black!60, anchor=south] at (10.0, 0.9) {Ambiente Emulado (QEMU)};

    \node[box] (out_shell)  at (0, 0)    {Shell\\(Gera Saída)};
    \node[box] (out_driver) at (4.0, 0)  {Driver UART\\(Escrita)};
    \node[box] (out_uart)   at (8.0, 0)  {UART PL011\\(Hardware)};
    \node[box] (out_term)   at (12.0, 0) {Terminal Serial\\(Emulador)};
    \node[box] (out_user)   at (16.0, 0) {Monitor\\(Visualização)};

    \draw[arrow] (out_shell) -- node[above=2pt, font=\scriptsize, text=black!80] {\texttt{printf()}} (out_driver);
    \draw[arrow] (out_driver) -- (out_uart);
    \draw[arrow] (out_uart) -- (out_term);
    \draw[arrow] (out_term) -- (out_user);

    \end{tikzpicture}%
    }
    \caption{Pipeline do fluxo de saída.}
    \label{fig:fluxo-out}
\end{figure}

\subsection{A Entrada de Dados via Teclado}

O teclado é um dispositivo de caractere. Diferentemente dos discos, que operam lendo e escrevendo dados em blocos de tamanho fixo, o teclado gera um fluxo sequencial e contínuo de caracteres independentes.

O driver implementado para gerenciar esse fluxo foi criado com base na especificação técnica da interface PL011 \cite{jankovic:qemu-uart-interface}. A comunicação em baixo nível é realizada mapeando na memória os principais registradores da controladora \cite{arm:summary-registers}, conforme detalhado na Tabela \ref{tab:uart-registradores}.


\begin{table}[htbp]
    \centering
    \renewcommand{\arraystretch}{1.3}
    \begin{tabular}{@{}llp{8cm}@{}}
        \hline
        \textbf{Registrador} & \textbf{Nome Completo} & \textbf{Descrição Funcional} \\
        \hline
        \texttt{UART0\_CR} & \textit{Control Register} & Habilita ou desabilita a UART, bem como as vias de transmissão (TX) e recepção (RX). \\
        \texttt{UART\_DR} & \textit{Data Register} & O \textit{kernel} lê deste endereço para obter dados recebidos e escreve nele para transmitir novos dados. \\
        \texttt{UART\_FR} & \textit{Flag Register} & Reflete o estado atual. Contém a flag \textbf{\texttt{RXFE}} (\textit{Receive FIFO Empty}), que indica se a fila de recepção está vazia. \\
        \texttt{UART0\_LCR\_H} & \textit{Line Control Register} & Controle de linha. Contém o bit \textbf{\texttt{FEN}} (\textit{FIFO Enable}) para habilitar as FIFOs de hardware. \\
        \texttt{UART0\_IMSC} & \textit{Interrupt Mask} & Define eventos geradores de interrupção. Destaque para as máscaras \textbf{\texttt{RXIM}} (recepção) e \textbf{\texttt{TXIM}} (transmissão). \\
        \texttt{UART0\_ICR} & \textit{Interrupt Clear} & Registrador de escrita; informa ao hardware que uma interrupção pendente foi devidamente tratada. \\
        \texttt{UART0\_MIS} & \textit{Masked Int. Status} & Registrador de leitura; indica exatamente quais interrupções (não mascaradas) estão pendentes no momento. \\
        \texttt{IBRD} e \texttt{FBRD} & \textit{Baud Rate Divisor} & Contêm, respectivamente, as partes inteira e fracionária do divisor para configurar a taxa de transmissão. \\
        \hline
    \end{tabular}
    \caption{Principais registradores da controladora UART PL011 utilizados na comunicação de E/S.}
    \label{tab:uart-registradores}
\end{table}

O \textit{kernel} efetiva o controle da UART lendo o \textit{status} e escrevendo configurações e dados diretamente nestes endereços mapeados.

\subsection{Inicialização da UART}
A rotina de inicialização ocorre por meio da chamada \texttt{uart\_init()}, que executa a seguinte sequência de passos, ilustrada na Figura \ref{fig:uart-init}.

\begin{figure}[htbp]
    \centering
    \begin{tikzpicture}[
        >=stealth,
        box/.style={
            rectangle,
            draw=black!70,
            thick,
            rounded corners=3pt,
            text width=7.1cm,
            align=center,
            fill=white,
            font=\small\sffamily,
            inner sep=6pt
        },
        arrow/.style={->, thick, draw=black!70}
    ]

    \node[box] (s1) at (0, 0)    {Desabilita temporariamente a UART\\\textit{(reconfiguração segura)}};
    \node[box] (s2) at (0, -2.1) {Ajusta a taxa de transmissão (\textit{baud rate})\\\texttt{uart\_set\_baud\_rate()}};
    \node[box] (s3) at (0, -4.2) {Ativa o mecanismo FIFO de hardware\\\texttt{uart\_enable\_fifo()}};
    \node[box] (s4) at (0, -6.3) {Mascara (habilita) as interrupções de recepção\\\texttt{uart\_enable\_interrupts()}};
    \node[box] (s5) at (0, -8.4) {Reabilita a UART para os fluxos de RX e TX};

    \draw[arrow] (s1) -- (s2);
    \draw[arrow] (s2) -- (s3);
    \draw[arrow] (s3) -- (s4);
    \draw[arrow] (s4) -- (s5);

    \end{tikzpicture}
    \caption{Sequência de inicialização da UART executada por \texttt{uart\_init()}.}
    \label{fig:uart-init}
\end{figure}

Ao final deste processo, o sistema encontra-se devidamente configurado para operar de forma assíncrona. Ao dispensar a espera ocupada (\textit{polling}), o \textit{kernel} delega ao controlador de interrupções a tarefa de notificá-lo apenas quando houver, de fato, novos caracteres a serem processados.
\subsection{Tratamento de Interrupções}

Sempre que um caractere é inserido no terminal do QEMU, a UART atualiza seu estado interno e levanta um sinal de interrupção (\textit{IRQ}). O controlador de interrupções da arquitetura intercepta esse sinal e desvia o fluxo de execução para a rotina de tratamento \texttt{uart\_irq\_handler(void)}, que realiza os seguintes passos lógicos:

\begin{enumerate}
    \item Verifica e valida a origem da interrupção lendo o registrador \texttt{UART0\_MIS};
    \item Confirma o tratamento do evento escrevendo no registrador de limpeza \texttt{UART0\_ICR};
    \item Consome os caracteres acumulados lendo diretamente o registrador \texttt{UART\_DR};
    \item Armazena em memória cada caractere recuperado chamando \texttt{uart\_buffer\_push()}, em um laço que se repete enquanto o hardware indicar que há dados válidos pendentes de leitura.
\end{enumerate}

\subsection{Gerenciamento de \textit{Buffer} Circular}

Como as velocidades de digitação e de processamento do \textit{kernel} são assimétricas, é necessário o armazenamento temporário dos dados recebidos. Para isso, estruturou-se uma fila circular, definida pelo tipo \texttt{ring\_buffer\_t}, cujos campos estão descritos na Tabela \ref{tab:ring-buffer-campos}.

\begin{table}[htbp]
    \centering
    \renewcommand{\arraystretch}{1.3}
    \begin{tabular}{@{}llp{7.5cm}@{}}
        \hline
        \textbf{Campo} & \textbf{Tipo} & \textbf{Descrição} \\
        \hline
        \texttt{buffer} & \texttt{char[]} & Vetor estático de caracteres em memória. \\
        \texttt{head} & \texttt{volatile uint32\_t} & Índice que rastreia a posição da próxima escrita. \\
        \texttt{tail} & \texttt{volatile uint32\_t} & Índice que rastreia a posição da próxima leitura. \\
        \hline
    \end{tabular}
    \caption{Campos da estrutura \texttt{ring\_buffer\_t}.}
    \label{tab:ring-buffer-campos}
\end{table}

A Figura \ref{fig:ring-buffer} ilustra o funcionamento do \textit{buffer} circular:

\begin{figure}[htbp]
    \centering
    \begin{tikzpicture}[
        >=stealth,
        cell/.style={
            rectangle,
            draw=black!70,
            thick,
            minimum width=1cm,
            minimum height=1cm,
            font=\scriptsize\sffamily
        },
        ptr/.style={
            ->,
            thick,
            draw=black!70
        }
    ]

    \foreach \i in {0,...,7} {
        \node[cell] (c\i) at (\i*1.1, 0) {};
    }

    \node[font=\scriptsize\sffamily] at (2*1.1, 1.3) {\texttt{tail}};
    \draw[ptr] (2*1.1, 1.0) -- (c2.north);

    \node[font=\scriptsize\sffamily] at (5*1.1, 1.3) {\texttt{head}};
    \draw[ptr] (5*1.1, 1.0) -- (c5.north);

    \draw[ptr, dashed] (c7.south) .. controls (4*1.1, -1.2) and (-1.2, -1.2) .. (c0.south);
    \node[font=\itshape\scriptsize, text=black!60] at (4*1.1, -1.5) {retorno ao início};

    \end{tikzpicture}
    \caption{Os índices \texttt{head} (escrita) e \texttt{tail} (leitura) avançam de forma independente e retornam ao início do vetor ao atingir a última posição.}
    \label{fig:ring-buffer}
\end{figure}

A partir dessa estrutura central, instancia-se o objeto \texttt{rx\_buffer} (dedicado à recepção), cujas operações de inserção e remoção ocorrem estritamente por meio da interface de funções do \textit{driver}:

\begin{description}
    \item[\texttt{static int uart\_buffer\_push(char c)}] insere o caractere \texttt{c} na posição \texttt{head} e avança o índice circularmente. Retorna \texttt{1} em caso de sucesso (havendo espaço livre) e \texttt{0} caso contrário;
    \item[\texttt{static int uart\_buffer\_pop(char *ptr)}] consome um caractere da posição \texttt{tail}, gravando-o na região de memória apontada por \texttt{ptr}. Retorna \texttt{1} se houver dado disponível e \texttt{0} se o \textit{buffer} estiver vazio.
\end{description}

\subsection{A Saída de Dados no Monitor}

Assim como o teclado, o monitor de texto é gerido pelo sistema como um dispositivo de caractere. No ecossistema emulado deste projeto, as saídas textuais fluem pela mesma controladora UART (PL011). Toda vez que o \textit{kernel} escreve um caractere para exibição, o hardware transmite essa informação de volta pela interface serial, encarregando o QEMU de renderizar os caracteres no terminal de saída visual.

%-------------------------------------------------------------------------------
\section{Interface Pública (API)}
%-------------------------------------------------------------------------------

O subsistema de E/S disponibiliza duas \textit{APIs}: uma de baixo nível, voltada ao acoplamento direto com o hardware, e outra de alto nível, estruturada para abstrair operações comuns ao restante do \textit{kernel}.

\subsection{API da UART (Baixo Nível)}

Esta camada agrupa as funções primitivas de interação direta com a controladora UART PL011. Suas principais rotinas são:

\begin{description}
    \item[\texttt{void uart\_init(void)}]
    Inicializa e configura a UART PL011 para comunicação serial. Durante o processo, a controladora é temporariamente desabilitada, a taxa de transmissão (\textit{baud rate}) é ajustada, os \textit{buffers} FIFO são ativados e as interrupções de recepção são habilitadas. Ao término, a UART é reativada para operações de transmissão e recepção.

    \item[\texttt{void uart\_putc(char c)}]
    Transmite um único caractere \texttt{c} pela UART. Realiza automaticamente o tratamento do caractere de nova linha (\texttt{\textbackslash n}), convertendo-o para a sequência (\texttt{\textbackslash r\textbackslash n}) exigida por terminais seriais. Antes de efetuar a escrita no registrador de dados \texttt{UART\_DR}, a função aguarda a liberação de espaço no FIFO de transmissão.

    \item[\texttt{void uart\_puts(const char *str)}]
    Transmite uma sequência de caracteres apontada pelo parâmetro \texttt{str}, iterando sobre o vetor e enviando cada caractere sucessivamente por meio de \texttt{uart\_putc()}. Ao encontrar o caracter \textbackslash n, finaliza a transmissão.

    \item[\texttt{char uart\_getc(void)}]
    Função sem parâmetros que retorna um caractere que foi recebido e transmitido pela UART. Ela funciona como uma interface pública de leitura do que a UART transmite.

    \item[\texttt{void uart\_irq\_handler(void)}]
    Deve ser acionada diretamente pelo controlador de interrupções da arquitetura (GIC) sempre que o identificador de interrupção correspondente (\textit{IRQ ID} 33) for disparado pelo hardware da UART. Inicia uma rotina de tratamentos de interrupções de UART.
\end{description}

\subsection{API de Entrada e Saída do Kernel (\texttt{kstdio})}

Implementada sobre o módulo \texttt{kstdio.c}, esta interface fornece abstrações de mais alto nível para a comunicação textual e o processamento de comandos do sistema, encapsulando as complexidades da controladora UART. Suas funções compreendem:

\begin{description}
    \item[\texttt{void kstdio\_init(void)}]
    Inicializa o subsistema básico de entrada e saída do \textit{kernel}, acionando internamente a rotina de configuração da UART.

    \item[\texttt{void kputs(const char *str)}]
    Envia uma \textit{string} completa (recebida através de um ponteiro) para a saída serial, servindo como interface de alto nível para exibição de mensagens estáticas do \textit{kernel}.

    \item[\texttt{void kputc(char c)}]
    Envia um único caractere para a saída serial utilizando as rotinas de abstração de E/S.

    \item[\texttt{char kgetc(void)}]
    Função sem parâmetros que retorna um caractere que foi recebido e transmitido pela UART. Ela funciona como uma interface pública de leitura do que a UART transmite.

    \item[\texttt{char kgets(char *buf, int max\_len)}]
    Lê uma sequência de caracteres da entrada padrão até encontrar uma quebra de linha ou atingir o limite de comprimento especificado em \texttt{max\_len}, armazenando o resultado no vetor de destino \texttt{buf}.

    \item[\texttt{void kprintf(const char *format, ...)}]
    Implementação personalizada de formatação e impressão para o \textit{kernel}, inspirada na clássica função \texttt{printf()} da biblioteca padrão de C. Interpreta uma \textit{string} de formatação contendo especificadores, convertendo argumentos subsequentes para representação textual e direcionando-os à saída serial. Os especificadores de formato suportados estão listados na Tabela \ref{tab:kprintf-especificadores}.
\end{description}

\begin{table}[htbp]
    \centering
    \renewcommand{\arraystretch}{1.2}
    \begin{tabular}{@{}lp{9cm}@{}}
        \hline
        \textbf{Especificador} & \textbf{Descrição} \\
        \hline
        \texttt{\%c} & Caractere isolado \\
        \texttt{\%s} & Cadeia de caracteres (\textit{string}) \\
        \texttt{\%d} & Número inteiro decimal com sinal \\
        \texttt{\%u} & Número inteiro decimal sem sinal \\
        \texttt{\%x} & Representação em formato hexadecimal \\
        \texttt{\%\textit{n}f} & Número de ponto flutuante, formatado com \textit{n} casas decimais \\
        \texttt{\%\%} & Caractere literal de porcentagem (\%) \\
        \hline
    \end{tabular}
    \caption{Especificadores de formato suportados por \texttt{kprintf()}.}
    \label{tab:kprintf-especificadores}
\end{table}


%-------------------------------------------------------------------------------
\section{Testes e Verificação}
%-------------------------------------------------------------------------------

Para validar o funcionamento integrado de todo o subsistema de E/S, utilizou-se o interpretador de comandos (\textit{shell}) implementado no arquivo \texttt{shell.c}. O \textit{shell} provê uma interface de linha de comando simples para interação direta com o \textit{kernel} através da porta serial.

O fluxo de execução do \textit{shell} baseia-se em um laço contínuo que exibe um indicador visual (\textit{prompt} \texttt{\$}) e aguarda a entrada do usuário por meio da função \texttt{kgetc()}. Os caracteres digitados são armazenados temporariamente em um \textit{buffer} interno de comandos. Durante essa captura, o sistema trata dois caracteres especiais de controle, descritos na Tabela \ref{tab:shell-teclas}.

\begin{table}[htbp]
    \centering
    \renewcommand{\arraystretch}{1.2}
    \begin{tabular}{@{}lp{9cm}@{}}
        \hline
        \textbf{Tecla} & \textbf{Efeito} \\
        \hline
        \textit{Enter} & Finaliza a digitação da linha atual, sinalizando que o comando está pronto para o processamento. \\
        \textit{Backspace} & Remove o último caractere inserido, atualizando visualmente o terminal. \\
        \hline
    \end{tabular}
    \caption{Tratamento de teclas especiais pelo \textit{shell}.}
    \label{tab:shell-teclas}
\end{table}

Para garantir a interatividade, cada caractere imprimível digitado é imediatamente enviado de volta ao terminal através de \texttt{kputc()}, gerando o efeito de eco visual. Após a conclusão da linha, o \textit{shell} executa uma etapa de \textit{parser} (tokenização), que divide a entrada com base em espaços, preenchendo o vetor de argumentos \texttt{argv} e atualizando a contagem em \texttt{argc}.

O reconhecimento e a execução dos comandos são efetuados por meio de comparações de \textit{strings} utilizando a função \texttt{strcmp()}. O conjunto de comandos implementados, descrito na Tabela \ref{tab:shell-comandos}, permitiu testar exaustivamente as rotinas do sistema.

\begin{table}[htbp]
    \centering
    \renewcommand{\arraystretch}{1.3}
    \begin{tabular}{@{}lp{9cm}@{}}
        \hline
        \textbf{Comando} & \textbf{Função e Cobertura de Teste} \\
        \hline
        \texttt{help} & Exibe a listagem de comandos disponíveis, testando a capacidade de impressão de múltiplas linhas formatadas via \texttt{kprintf()}. \\
        \texttt{clear} & Limpa o terminal serial enviando a sequência de escape ANSI padrão (\texttt{\textbackslash033[2J\textbackslash033[H}) por meio da função \texttt{kputs()}. \\
        \texttt{echo} & Reproduz no terminal os argumentos subsequentes informados pelo usuário. Valida a leitura integrada do \textit{buffer}, a tokenização de \textit{strings} e a transmissão formatada de texto. \\
        \texttt{kprintf} & Executa uma bateria de testes da função de impressão formatada, validando a conversão e exibição de múltiplos especificadores (\texttt{\%s}, \texttt{\%c}, \texttt{\%d} e \texttt{\%x}). \\
        \hline
    \end{tabular}
    \caption{Comandos do \textit{shell} utilizados na verificação do subsistema de E/S.}
    \label{tab:shell-comandos}
\end{table} 
\chapter{Chamadas de Sistema}
\chapterauthor{Cauã Marques da Silva, Felipe de Castro Duchen Auroux, Felipe de Oliveira Renó, Vicente Pinto Tomás Junior}
%-------------------------------------------------------------------------------
\section{Descrição}
%-------------------------------------------------------------------------------

As chamadas de sistema são o mecanismo pelo qual o sistema operacional fornece aos processos acesso aos recursos do sistema. Essencialmente, as chamadas de sistema são a forma pela qual os processos têm acesso aos arquivos, diretórios, ao relógio, teclado e até mesmo a outros processos.

O sistema operacional oferece esse mecanismo para poder manter a separação entre os dois principais modos de operação do processador: supervisor e usuário. No modo supervisor o processador pode executar qualquer instrução e interagir com outros componentes físicos do sistema; já no modo usuário o processador está restrito a apenas um subconjunto dessas instruções e não pode interagir com outros componentes físicos do sistema. Esse isolamento busca impedir que um processo, defeituoso ou malicioso, assuma controle dos recursos do sistema.

Além disso, a interface criada por esse mecanismo oferece aos programas uma abstração dos recursos do sistema, ou seja, um programa não precisa conhecer o hardware responsável pela memória se ele deseja salvar um arquivo, basta que ele peça ao sistema operacional. Ademais, o programa não precisa nem saber como a informação será efetivamente salva na memória, ele apenas precisa fornecer o que deseja salvar e o sistema operacional lidará com o resto.

%-------------------------------------------------------------------------------
\section{Decisões de Arquitetura e Design}
%-------------------------------------------------------------------------------

As chamadas de sistema estão diretamente ligadas à forma como o núcleo é estruturado; na literatura discutem-se, principalmente, duas arquiteturas: núcleo monolítico e micronúcleo.

A arquitetura monolítica, essencialmente, se baseia em concentrar as operações de nível de sistema dentro do núcleo, ou seja, a gerência de memória, o sistema de arquivos, as pilhas dos protocolos de rede e outros existem no espaço do núcleo (são executados em modo supervisor).

Na arquitetura micronúcleo as operações de nível de sistema são transferidas para o espaço do usuário (são executadas em modo usuário) por módulos próprios separados do núcleo. O núcleo fica, portanto, responsável por escalonar e gerenciar a comunicação entre os processos que agora, também, incluem os módulos do sistema.

Dessa forma, as chamadas de sistema em um núcleo monolítico são executadas em modo supervisor e interagem diretamente com os recursos do sistema; em um micronúcleo, as chamadas seriam transformadas em trocas de mensagens entre o processo fazendo a chamada e o processo responsável pelo recurso do sistema, de forma que o micronúcleo funcionaria apenas como um mediador encaminhando essas mensagens. A Figura \ref{fig:arquiteturas-nucleo} ilustra as principais diferenças entre as duas abordagens.

\begin{figure}[htbp]
    \centering
    \resizebox{\textwidth}{!}{%
    \begin{tikzpicture}[
        font=\sffamily,
        layer/.style={draw=black, thick, minimum height=1cm, align=center, inner sep=4pt},
    ]

    %====================== MICRONÚCLEO (esquerda) — topo em y=6.6 ======================
    \node[layer, fill=black!85, text=white, minimum width=7cm] (mk-hw) at (3.5, 0.5) {\Large Hardware};

    \node[layer, fill=white, minimum width=7cm, minimum height=2.0cm] (mk-kernel) at (3.5, 2.2) {};
    \node[font=\sffamily, anchor=north west] at (0.2, 3.05) {IPC Básico};
    \node[layer, fill=white, minimum width=3.1cm, minimum height=0.5cm] (mk-mm) at (5.0, 2.85) {\footnotesize Gerência de Memória};
    \node[font=\sffamily\bfseries\Large] at (3.5, 1.75) {Escalonamento};

    \node[layer, fill=black!55, text=white, minimum width=1.72cm, minimum height=1.9cm, text width=1.55cm, font=\sffamily\footnotesize] (mk-ipc)   at (0.86, 4.35) {IPC de Aplicação};
    \node[layer, fill=black!30, minimum width=1.72cm, minimum height=1.9cm, text width=1.55cm, font=\sffamily\footnotesize]             (mk-proto) at (2.61, 4.35) {Pilha de Protocolos};
    \node[layer, fill=black!70, text=white, minimum width=1.72cm, minimum height=1.9cm, text width=1.55cm, font=\sffamily\footnotesize] (mk-drv)   at (4.36, 4.35) {Driver de Dispositivo};
    \node[layer, fill=black!12, minimum width=1.72cm, minimum height=1.9cm, text width=1.55cm, font=\sffamily\footnotesize]             (mk-fs)    at (6.11, 4.35) {Servidor de Arquivos};

    \node[layer, fill=black!45, text=white, minimum width=7cm] (mk-app) at (3.5, 6.1) {\Large Aplicação};

    \foreach \x in {0.86, 2.61, 4.36, 6.11} {
        \draw[<->, thick] (\x, 5.35) -- (\x, 5.55);
    }

    \node[align=center, font=\bfseries\Large] at (3.5, -0.55) {Micronúcleo};

    %====================== SEPARADOR VS ======================
    \fill[black!15] (7.3,1.2) -- (10.7,1.2) -- (10.7,5.5) -- (7.3,3.2) -- cycle;
    \node[font=\bfseries\Huge] at (9.0, 3.45) {Vs};
    \node[font=\sffamily, align=center] at (9.0, 5.9) {Espaço de\\Usuário};
    \node[font=\sffamily, align=center] at (9.0, 1.9) {Espaço de\\Núcleo};

    %====================== NÚCLEO MONOLÍTICO (direita) — topo em y=6.6 ======================
    \node[layer, fill=black!85, text=white, minimum width=7cm] (mono-hw) at (14.5, 0.5) {\Large Hardware};

    \node[layer, fill=white, minimum width=7cm, minimum height=0.9cm] (mono-drv) at (14.5, 1.57) {\large Driver de Dispositivo};

    \node[layer, fill=white, minimum width=7cm, minimum height=1.3cm] (mono-sched) at (14.5, 2.79) {};
    \node[layer, fill=white, minimum width=3.1cm, minimum height=0.5cm] (mono-mm) at (16.0, 3.19) {\footnotesize Gerência de Memória};
    \node[font=\sffamily\bfseries\large] at (14.5, 2.55) {Escalonador};

    \node[layer, fill=white, minimum width=7cm, minimum height=0.9cm] (mono-ipc) at (14.5, 4.01) {\large IPC, Sistema de Arquivos};

    \node[layer, fill=white, minimum width=7cm, minimum height=0.9cm] (mono-sc) at (14.5, 5.03) {\large Chamadas de Sistema};

    \node[layer, fill=black!45, text=white, minimum width=7cm] (mono-app) at (14.5, 6.1) {\Large Aplicação};

    \node[align=center, font=\bfseries\Large] at (14.5, -0.55) {Núcleo Monolítico};

    \end{tikzpicture}%
    }
    \caption{Comparação entre as arquiteturas de micronúcleo e núcleo monolítico.}
    \label{fig:arquiteturas-nucleo}
\end{figure}

Nesse contexto, para o desenvolvimento das chamadas de sistema neste projeto, adotou-se uma abordagem de núcleo monolítico visando simplificar os detalhes de implementação das chamadas em si, uma vez que as mesmas podem interagir diretamente com os recursos do sistema. Consequentemente, essa abordagem acarreta em um núcleo mais complexo (no que se refere à quantidade de partes e responsabilidades do núcleo), o que reduz a modularidade e introduz mais pontos de falha no sistema (uma vez que as chamadas de sistema acontecem em modo supervisor, o que pode ser perigoso caso haja falhas no código).

%-------------------------------------------------------------------------------
\section{Detalhes da Implementação}
%-------------------------------------------------------------------------------

Há quatro principais componentes envolvidos em realizar uma chamada de sistema, sendo eles: o invocador, o tratador SVC, o tratador de chamada de sistema e a tabela de chamadas de sistema. A Figura \ref{fig:syscall-flow} ilustra o fluxo de execução entre esses quatro componentes, detalhado nos parágrafos a seguir.

\begin{figure}[htbp]
    \centering
    \begin{tikzpicture}[
        >=stealth,
        box/.style={
            rectangle,
            draw=black!70,
            thick,
            rounded corners=3pt,
            text width=9.2cm,
            align=left,
            fill=white,
            font=\small\sffamily,
            inner sep=6pt
        },
        arrow/.style={->, thick, draw=black!70}
    ]

    \node[box] (b1) at (0, 0)     {\textbf{Invocador} (rotina em C)\\Salva o ponto de retorno, define a chamada desejada e executa \texttt{SVC \#0}};
    \node[box] (b2) at (0, -2.3)  {\textbf{Tratador SVC} --- \texttt{svc\_handler}\\Salva os registradores comuns e transfere a execução para \texttt{syscall\_Handler}};
    \node[box] (b3) at (0, -4.6)  {\textbf{Tratador de Chamada de Sistema} --- \texttt{syscall\_Handler}\\Consulta a \texttt{Syscall\_Table} usando o identificador da chamada (\texttt{r0})};
    \node[box] (b4) at (0, -6.9)  {\textbf{Tabela de Chamadas de Sistema} --- \texttt{Syscall\_Table}\\Aponta para o endereço da rotina correspondente (ex.: \texttt{printf\_handler})};
    \node[box] (b5) at (0, -9.2)  {\textbf{Rotina da Chamada de Sistema}\\Executa a operação solicitada e devolve o resultado ao processo invocador};

    \draw[arrow] (b1) -- (b2);
    \draw[arrow] (b2) -- (b3);
    \draw[arrow] (b3) -- (b4);
    \draw[arrow] (b4) -- (b5);

    \end{tikzpicture}
    \caption{Fluxo de execução de uma chamada de sistema, através dos quatro componentes envolvidos.}
    \label{fig:syscall-flow}
\end{figure}

Os invocadores são as rotinas encapsuladas pelas funções em C que oferecem as chamadas de sistema para os processos. Essas precisam registrar o ponto atual da execução (para que se possa retornar a esse ponto posteriormente), registrar o endereço da rotina que será executada e, então, gerar uma exceção forçando a transferência da execução para o sistema operacional em modo supervisor, como no exemplo do Código \ref{lst:invocador}.

\begin{lstlisting}[language=C, caption={Exemplo de invocador}, label={lst:invocador}]
exemplo:
    push {lr}
    mov r0, #CHAMADA
    SVC #0 
    pop {pc}
\end{lstlisting}

Quando a instrução SVC é usada, o processador gera uma exceção; ele então consulta a tabela de exceções e transfere a execução para a rotina identificada na tabela. O tratador SVC, simplesmente, salva os registradores comuns e transfere a execução para o tratador de chamada de sistemas (Código \ref{lst:svc-handler}).

\begin{lstlisting}[language=C, caption={Tratador SVC}, label={lst:svc-handler}]
svc_handler:
    STMFD sp!, {r1-r12, lr}
    bl syscall_Handler
    LDMFD sp!, {r1-r12, pc}^
\end{lstlisting}

O tratador de chamadas de sistema interpretará pelos argumentos da chamada qual é a intenção do processo e fará uma consulta na tabela de chamadas de sistema (Código \ref{lst:syscall-handler}).

\begin{lstlisting}[language=C, caption={Tratador de chamada de sistema}, label={lst:syscall-handler}]
syscall_Handler:
    ldr r3, =Syscall_Table
    ldr r3, [r3,r0,LSL #2]
    bx r3
\end{lstlisting}

Se possível, ele executará a rotina associada ao endereço na tabela (Código \ref{lst:syscall-table}) e devolverá o resultado e a execução para o processo que fez a chamada de sistema.

\begin{lstlisting}[language=C, caption={Tabela dos endereços das chamadas de sistema}, label={lst:syscall-table}]
Syscall_Table:
    .word printf_handler
    .word abort_handler
\end{lstlisting}

\begin{lstlisting}[language=C, caption={Exemplo de rotina que será de fato executada}, label={lst:printf-handler}]
printf_handler:
    mov r0,r1
    push {lr}
    bl serial_puts
    mov r0, #10
    pop {pc}
\end{lstlisting}


%-------------------------------------------------------------------------------
\section{Interface Pública (API)}
%-------------------------------------------------------------------------------

O subsistema de chamadas de sistema disponibiliza, por ora, duas chamadas, descritas na Tabela \ref{tab:api-syscalls}.

\begin{table}[htbp]
    \centering
    \renewcommand{\arraystretch}{1.3}
    \newcolumntype{Y}{>{\raggedright\arraybackslash}X}
    \begin{tabularx}{\textwidth}{@{}l >{\hsize=0.85\hsize}Y >{\hsize=0.75\hsize}Y >{\hsize=1.4\hsize}Y@{}}
        \hline
        \textbf{Função} & \textbf{Parâmetros} & \textbf{Retorno} & \textbf{Descrição} \\
        \hline
        \texttt{printf} & Ponteiro de \texttt{char} para o texto a ser impresso & \texttt{0} (sucesso)\newline\texttt{-1} (falha) & Imprime o texto passado como parâmetro na saída serial. \\
        \texttt{abort} & --- & \texttt{0} (sucesso)\newline\texttt{-1} (falha)$^{*}$ & Força o reinício da CPU. Em processadores ARM é necessário interagir com o PSCI (\textit{Power State Coordination Interface}). \\
        \hline
    \end{tabularx}
    \caption{Chamadas de sistema disponíveis na API pública.}
    \label{tab:api-syscalls}
\end{table}

{\footnotesize $^{*}$Observação: por enquanto a chamada \texttt{abort} não possui uma forma de retornar nenhum valor, uma vez que o reinício é instantâneo.}
%-------------------------------------------------------------------------------
\section{Testes e Verificação}
%-------------------------------------------------------------------------------

As rotinas que invocam as chamadas de sistema estão encapsuladas em C, então os testes consistiram em invocá-las, verificar, pelo depurador, as alterações no fluxo de execução e observar o resultado da execução das rotinas. A Tabela \ref{tab:testes-syscalls} resume os resultados obtidos.

\begin{table}[htbp]
    \centering
    \small
    \renewcommand{\arraystretch}{1.3}
    \begin{tabular}{@{}p{1.7cm}p{4.6cm}p{4.6cm}@{}}
        \hline
        \textbf{Chamada} & \textbf{Comportamento Esperado} & \textbf{Comportamento Obtido} \\
        \hline
        \texttt{printf} & Mensagem ``Testando alterações'' exibida na tela & Mensagem ``Testando alterações'' exibida na tela \\
        \texttt{abort} & Reinício da CPU, com a mensagem ``Executando em modo ARM bare-metal no QEMU.'' exibida na tela & Salto do endereço corrente para o endereço \textit{start} (sinalizando o reinício), seguido da exibição da mensagem ``Executando em modo ARM bare-metal no QEMU.'' \\
        \hline
    \end{tabular}
    \caption{Resultados dos testes das chamadas de sistema.}
    \label{tab:testes-syscalls}
\end{table}

Essas chamadas de sistema apenas encapsulam rotinas que já existiam (no caso do \texttt{printf}) ou instruções especiais da CPU (no caso do \texttt{abort}); portanto, os casos de borda se restringiram ao que acontece durante uma chamada de sistema, ou seja, se os registradores estavam sendo manipulados corretamente (salvos e depois restaurados), o que foi constatado pela inspeção durante a depuração de ambas as chamadas.
\include{capitulos/02_gerenciamento_de_memoria}  
\chapter{Processos, Troca de Contexto e Escalonamento}

\chapterauthor{João Pedro Rizzo Sales, Samuel Romano Ribeiro, Guilherme Michelsen do Monte}

\chapterrev{Ti Ribeiro Silvério}


\section{Descrição}


Este capítulo documenta três módulos do UFSKernel:
o módulo de \textbf{processos}, responsável por representar, criar e encerrar
unidades de execução; o mecanismo de \textbf{troca de contexto}, que realiza
a alternância segura entre modo usuário e modo kernel; e o \textbf{escalonador},
que decide qual processo recebe a CPU e por quanto tempo. Juntos, esses módulos
viabilizam a multiprogramação no kernel, conforme descrito em
Silberschatz et al.~\cite{silberschatz2018} e Tanenbaum e Bos~\cite{tanenbaum2015}.

Os arquivos que compõem esses módulos são: \texttt{process.h},
\texttt{process.c}, \texttt{proc\_ll.c}, \path{svc_entry.S}, \texttt{irq\_entry.S}, \texttt{irq\_return.S}, \texttt{context\_switch.S}, \texttt{scheduler.h} e \path{process_control_block.c}.


\section{Fundamentação Teórica}


\subsection{Processo e PCB}

Um processo é um programa em execução: uma entidade dinâmica que engloba
o contador de programa, os registradores, a pilha e o espaço de endereços
ocupado na memória~\cite{tanenbaum2015}. Para gerenciá-los, o kernel mantém
um \textbf{Bloco de Controle de Processo} (PCB) por processo, concentrando
tudo o que é necessário para suspendê-lo e retomá-lo como se nunca tivesse
sido interrompido~\cite{silberschatz2018}.

Durante seu ciclo de vida, um processo transita entre quatro estados. Os três
canônicos, \textbf{Pronto}, \textbf{Em execução} e \textbf{Bloqueado}, são
descritos em~\cite{tanenbaum2015}. O UFSKernel acrescenta um quarto,
\textbf{Zumbi}: o processo encerrou, mas sua estrutura ainda não foi
descartada, pois o processo pai deve registrar o encerramento do filho antes
que os recursos sejam liberados~\cite{silberschatz2018}.

\subsection{Trapframe, Contexto de Kernel e ABI ARM}

Na ARMv7, o processador opera em modos de privilégio distintos. Toda transição
de modo usuário para kernel, seja por chamada de sistema (\texttt{SVC}) ou por
interrupção (\texttt{IRQ}), exige salvar o estado completo do processo para que
a execução possa ser retomada. A estrutura que armazena
esse estado é o \textbf{trapframe}~\cite{tanenbaum2015}.

Para trocas de contexto internas ao kernel, quando o escalonador substitui
um processo por outro sem sair do modo supervisor, apenas um subconjunto
dos registradores precisa ser salvo. A ABI ARM (AAPCS)~\cite{arm:armv7ar} define
dois grupos: registradores \emph{call-preserved} (\texttt{r4}--\texttt{r11}, SP e LR), que uma
função deve restaurar antes de retornar, e registradores \emph{call-clobbered}
(\texttt{r0}--\texttt{r3}, \texttt{r12}), que o chamador sabe que podem ser destruídos. Como o compilador
já protege os call-clobbered em torno de qualquer chamada de função, o
mecanismo de troca de contexto de kernel precisa salvar apenas os call-preserved.

\subsection{Escalonamento: Round-Robin e MLFQ}

O escalonador decide qual processo pronto recebe a CPU e por quanto tempo
(o \textbf{quantum}), tipicamente acionado por um timer periódico~\cite{silberschatz2018}.
O \textbf{Round-Robin} organiza os processos em fila circular com quantum fixo,
mas trata igualmente processos interativos e CPU-bound~\cite{tanenbaum2015}.
A \textbf{Multilevel Feedback Queue} (MLFQ) supera essa limitação: mantém
múltiplas filas de prioridade com quanta crescentes, promovendo processos que
usam pouco CPU (interativos) e rebaixando os que esgotam o quantum (CPU-bound),
de modo que cada processo migra naturalmente para o nível compatível com seu
comportamento real~\cite{silberschatz2018}.


\section{Módulo de Processos}


\subsection{Decisões de Arquitetura e Design}

O módulo mantém duas estruturas distintas para o estado de execução de um
processo. O \texttt{struct trapframe} armazena todos os registradores de
propósito geral (\texttt{r0}--\texttt{r12}), os registradores banked do modo usuário
(\texttt{sp\_usr}, \texttt{lr\_usr}), o contador de programa (\texttt{pc\_usr})
e o registrador de estado (\texttt{cpsr\_usr}) — tudo o que o processo de
usuário enxerga. O \texttt{struct context} armazena apenas os registradores
call-preserved (\texttt{r4}--\texttt{r11}, SP e LR), suficientes para a troca de contexto
dentro do kernel. Salvar \texttt{r0}--\texttt{r3} e \texttt{r12} no contexto de kernel seria redundante,
pois o compilador já os trata como voláteis em torno de qualquer
chamada~\cite{arm:armv7ar}.

\subsubsection{Alocação Dinâmica de Memória}

As estruturas
fundamentais de um processo — o próprio PCB (\texttt{struct process}),
a pilha de usuário (\texttt{usr\_stack}) e a pilha de kernel (\texttt{kstack}) —
são alocadas dinamicamente requisitando blocos de 4~KB ao PMM através da 
função \texttt{pmm\_alloc\_block()}. O limite de 64 processos simultâneos existe para simplificar a gestão de PIDs, mas a memória física só é
efetivamente consumida sob demanda.

\subsubsection{Execução de Programas e Passagem de Argumentos}

Sem um sistema de arquivos externo disponível, os executáveis são representados
por funções compiladas juntamente com o kernel, gerenciadas pela
\texttt{program\_table} (definida em \path{programas.h}).

\lstset{captionpos=t}
\begin{lstlisting}[language=C, caption={Assinatura padronizada dos programas}]
void (*programa)(int argc, char** argv);
\end{lstlisting}
\lstset{captionpos=b}

Essa abordagem permite testar o ciclo completo \texttt{fork}--\texttt{exec}--\texttt{exit} 
simulando a passagem real de argumentos (\texttt{argc} e \texttt{argv}) para o 
espaço de execução do processo de forma análoga a sistemas POSIX reais.

\subsubsection{Gestão de PIDs e Bloqueio de Processos}

PIDs são gerenciados por \texttt{create\_pid} e \texttt{free\_pid} sobre
o array \path{used_pids[MAX_PROCESS_COUNT]}. O PID~0 é
reservado, convenção Unix que evita confundi-lo com um ponteiro
nulo~\cite{silberschatz2018}. O sistema emprega ativamente o estado de
bloqueio (\texttt{BLOCKED}). Um processo pode ser bloqueado ao aguardar por 
entrada de dados (E/S via UART) ou ao sincronizar com o término de um filho 
através da chamada \texttt{waitpid}, sendo reinserido na fila de prontos
assim que a condição de espera for satisfeita~\cite{tanenbaum2015}.

\subsection{Detalhes da Implementação}

\subsubsection{O PCB: \texttt{struct process}}

\lstset{captionpos=t}
\begin{lstlisting}[language=C, caption={Bloco de Controle de Processo em \texttt{process.h}}]
struct process {
    char* usr_stack;        // inicio da pilha de usuario
    char* kstack;           // inicio da pilha de kernel
    struct trapframe *tf;   // ponteiro para o trapframe (offset 8)
    enum State state;       // RUNNING, READY, BLOCKED ou ZOMBIE
    enum Block_Reason blocked_by;
    pid_t waiting_for_pid;  // PID do filho aguardado, se houver
    pid_t pid;
    process parent;
    pid_list children_ids;
    struct context context; // registradores call-preserved (troca de contexto)
    unsigned priority;
};
\end{lstlisting}
\lstset{captionpos=b}

A ordem dos três primeiros campos é estritamente deliberada: \texttt{usr\_stack} 
(4~bytes) e \texttt{kstack} (4~bytes) posicionam \texttt{tf} exatamente
no offset~8 — valor codificado na constante \texttt{TRAPFRAME\_OFFSET = 8},
amplamente utilizada pelo código assembly para acessar o trapframe do processo
diretamente.

\subsubsection{Ciclo de vida: \texttt{first\_process}, \texttt{fork}, \texttt{exec}, \texttt{exit} e \texttt{waitpid}}

\texttt{first\_process} cria o processo inicial: aloca as estruturas de pilha via PMM,
insere o processo no escalonador no nível de maior prioridade, define
\texttt{current} e invoca \texttt{sys\_exec} passando o ponteiro da função
solicitada.

\texttt{sys\_fork} cria um clone do processo. Ele requisita novas páginas ao
PMM para o \texttt{usr\_stack} e \texttt{kstack} do filho e copia o conteúdo
byte a byte da memória do pai. Os ponteiros do trapframe são meticulosamente
recalculados por aritmética de ponteiros baseada no deslocamento entre as pilhas:

\lstset{captionpos=t}
\begin{lstlisting}[language=C, caption={Ajuste de deslocamento na pilha no \texttt{sys\_fork}}]
unsigned int stack_diff = 
    (unsigned int)newprocess->usr_stack - (unsigned int)newprocess->parent->usr_stack;
newprocess->tf->sp_usr = newprocess->parent->tf->sp_usr + stack_diff;

unsigned int *saved_r7 = (unsigned int *)newprocess->tf->sp_usr;
*saved_r7 = *saved_r7 + stack_diff;
\end{lstlisting}
\lstset{captionpos=b}

O pai recebe o PID do filho em \texttt{\texttt{r0}}; o filho recebe zero, convenção
Unix clássica~\cite{silberschatz2018}. O \texttt{context.lr} do filho aponta
para \texttt{fork\_return\_asm}, garantindo que ele desembarque adequadamente
na restauração do modo usuário ao ser escalonado pela primeira vez.

\texttt{sys\_exec} prepara a imagem para a nova execução. Ele configura o 
\texttt{pc\_usr} para o ponteiro da nova função e repassa \texttt{argc} (em \texttt{\texttt{r0}}) e 
\texttt{argv} (em \texttt{r1}). Por motivos de isolamento e depuração, os 
registradores restantes (\texttt{r2} a \texttt{\texttt{r12}}) são zerados explicitamente. Em seguida, 
invoca-se \texttt{restore\_user\_context}.

\texttt{sys\_exit} marca o estado do processo atual como \texttt{ZOMBIE}. Crucialmente, 
verifica se o seu \texttt{parent} encontra-se bloqueado em \texttt{waitpid} aguardando
seu término (\texttt{waiting\_for\_pid}). Se afirmativo, muda o estado do pai para 
\texttt{READY}, efetivamente acordando-o. Por fim, invoca \texttt{pcb\_elect} 
para ceder imediatamente a CPU.

\texttt{sys\_waitpid} é o mecanismo responsável pela coleta (\emph{reaping}) de
processos Zumbis. Ao ser invocada, bloqueia o pai até que o processo filho 
especificado termine (mude para \texttt{ZOMBIE}). Uma vez encerrado, 
\texttt{waitpid} libera os blocos físicos de memória do filho de volta 
ao PMM, descarta seu PCB da tabela global, libera o PID na estrutura circular 
e retoma a execução do pai.

\section{Mecanismo de Troca de Contexto}


\subsection{Três mecanismos: entrada, retorno e troca}

O mecanismo de troca de contexto é composto por três partes complementares:
a \textbf{entrada} (\texttt{svc\_entry.S}, \texttt{irq\_entry.S}), que captura
o estado do processo de usuário e o armazena em um trapframe na pilha de kernel;
o \textbf{retorno} (\texttt{irq\_return.S}), que restaura esse estado e devolve
o controle ao processo; e a \textbf{troca de contexto de kernel}
(\texttt{context\_switch.S}), que salva e restaura apenas os registradores
call-preserved para alternar entre processos enquanto permanece em modo supervisor.
Mantê-los separados permite ao escalonador decidir, a cada interrupção, se
retorna ao mesmo processo ou realiza uma troca.

\subsection{Entrada via SVC e IRQ}

Ao receber um SVC (System Call) ou IRQ (Interrupção de Hardware), o processador ARMv7
salva automaticamente o CPSR em SPSR e o endereço de retorno em LR~\cite{arm:armv7ar}. 
A construção do trapframe, no entanto, deve ocorrer obrigatoriamente na pilha 
de kernel do processo (\texttt{SP\_svc}). 

Para chamadas de sistema (\texttt{svc\_entry.S}), o processador já ingressa 
no modo Supervisor, permitindo alocar o espaço e salvar o contexto de imediato.
Porém, para interrupções (\texttt{irq\_entry.S}), o processador entra no modo IRQ,
que possui sua própria pilha (\texttt{SP\_irq}). A solução empregada no UFSKernel é
salvar temporariamente os registradores mais voláteis na pilha IRQ, transitar 
manualmente para o modo SVC alterando o CPSR, alocar o trapframe na kstack e 
retornar brevemente ao modo IRQ para resgatar os valores originais:

\lstset{captionpos=t}
\begin{lstlisting}[language={[x86masm]Assembler},
                   caption={Transição de modo e resgate de registradores em \texttt{irq\_entry.S}}]
stmfd sp!, {\texttt{r0}-\texttt{r3}, lr}       @ Salva \texttt{r0}-\texttt{r3} e LR na pilha IRQ
mrs \texttt{r0}, cpsr
bic \texttt{r0}, \texttt{r0}, #0x1F
orr \texttt{r0}, \texttt{r0}, #0x13
msr cpsr_c, \texttt{r0}               @ Troca para o modo SVC (0x13)

sub sp, sp, #(17 * 4)        @ Aloca o trapframe na kstack (SP_svc)
add r1, sp, #(4 * 4)
stmia r1, {\texttt{r4}-\texttt{r12}}           @ Salva \texttt{r4}-\texttt{r12} em seguranca

@ ... (Transita de volta ao IRQ para ejetar a pilha e joga o LR em \texttt{r4})
\end{lstlisting}
\lstset{captionpos=b}

Após estabilizar a execução no modo Supervisor com os dados protegidos, 
finaliza-se a construção do trapframe:

\lstset{captionpos=t}
\begin{lstlisting}[language={[x86masm]Assembler},
                   caption={Finalização do trapframe na pilha de kernel}]
stmia sp, {\texttt{r0}-\texttt{r3}}            @ Salva \texttt{r0}-\texttt{r3} no inicio do trapframe
str \texttt{r4}, [sp, #(15 * 4)]      @ Salva o PC (resgatado em \texttt{r4})
str r1, [sp, #(16 * 4)]      @ Salva o SPSR original

add r2, sp, #(13 * 4)
stmia r2, {sp, lr}^          @ Salva SP_usr e LR_usr (^ acessa banked USR)
\end{lstlisting}
\lstset{captionpos=b}

O sufixo \texttt{\^} em \texttt{stmia r2, \{sp, lr\}\^} é essencial: ele
instrui o processador a acessar os registradores banked do modo usuário
(\texttt{SP\_usr} e \texttt{LR\_usr}), e não os do modo corrente. Sem ele,
o estado do processo seria corrompido~\cite{arm:armv7ar}. Antes de chamar
código C, a pilha é alinhada em 8 bytes, respeitando a convenção AAPCS.

\subsection{Retorno ao Modo Usuário}

\texttt{irq\_return.S} realiza o processo inverso: lê o trapframe no topo da
kstack, restaura o CPSR do usuário no SPSR, carrega \texttt{r0}--\texttt{r12}, restaura
\texttt{sp\_usr} e \texttt{lr\_usr} com o sufixo \texttt{\^}, ejeta o 
espaço da pilha e executa o retorno:

\lstset{captionpos=t}
\begin{lstlisting}[language={[x86masm]Assembler},
                   caption={Retorno ao modo usuário em \texttt{irq\_return.S}}]
ldr \texttt{r0}, [sp, #(16 * 4)]      @ \texttt{r0} = SPSR salvo
msr spsr_cxsf, \texttt{r0}            @ joga para SPSR_svc

ldr lr, [sp, #(15 * 4)]      @ carrega o PC salvo para LR_svc
add \texttt{r0}, sp, #(13 * 4)
ldmia \texttt{r0}, {sp, lr}^          @ restaura SP_usr e LR_usr
nop                              
ldmia sp, {\texttt{r0}-\texttt{r12}}           @ restaura \texttt{r0}-\texttt{r12}
add sp, sp, #(17 * 4)        @ Ejetamos o trapframe da kstack!

movs pc, lr                  @ Pula para o PC e restaura o CPSR nativamente
\end{lstlisting}
\lstset{captionpos=b}

O sufixo \texttt{s} na instrução \texttt{movs} com o PC como destino em 
modo privilegiado copia automaticamente o SPSR para o CPSR~\cite{arm:armv7ar} — 
restaurando modo, flags de processador e conjunto de instruções em uma única 
operação atômica.

\subsection{Troca de Contexto de Kernel}

\lstset{captionpos=t}
\begin{lstlisting}[language={[x86masm]Assembler},
                   caption={Troca de contexto de kernel em \texttt{context\_switch.S}}]
context_switch:
    cmp \texttt{r0}, #0
    beq if_null              @ se nao ha processo saindo, apenas carrega

    stmia \texttt{r0}!, {\texttt{r4}-\texttt{r11}}      @ salva registradores call-preserved
    str sp, [\texttt{r0}], #4         @ salva SP
    str lr, [\texttt{r0}]             @ salva LR (ponto de retorno quando reativado)

if_null:
    ldmia r1!, {\texttt{r4}-\texttt{r11}}      @ restaura registradores call-preserved
    ldr sp, [r1], #4         @ restaura SP (agora na pilha do novo processo)
    ldr lr, [r1]             @ restaura LR

    bx lr                    @ retorna -- no contexto do processo entrante
\end{lstlisting}
\lstset{captionpos=b}

O caso \texttt{\texttt{r0} == NULL} cobre o boot inicial, quando não há um 
processo anterior saindo da CPU. A propriedade central do \texttt{context\_switch} 
é que o \texttt{bx lr} não retorna para quem chamou a função de troca, mas para 
o ponto exato em que o processo entrante havia sido suspenso. Para um processo 
recém-criado via \texttt{fork}, esse ponto de entrada é forjado como 
\texttt{fork\_return\_asm} (que redireciona para \texttt{irq\_return}), 
fazendo com que o novo fluxo de execução desembarque diretamente e de 
forma segura no modo usuário.

\section{Escalonador}

\subsection{Decisões de Arquitetura e Design}

\subsubsection{MLFQ com 32 níveis e quanta crescentes}

A ideia adotada é uma MLFQ simplificada com 32 níveis. O \texttt{pcb}
é um array de \texttt{struct pcb\_entry} com quanta variando de 20 a 175
em progressão aritmética de passo 5. Processos recém-criados entram no
nível~0, maior prioridade, menor quantum, refletindo a heurística de que
processos novos tendem a ser interativos~\cite{tanenbaum2015}. A progressão
aritmética produz separação gradual entre os níveis sem saltos abruptos que
pudessem causar \emph{starvation} nos níveis inferiores.

\subsubsection{Margem de quantum e \texttt{list\_location}}

A constante \texttt{SCHEDULER\_QUANTUM\_MARGIN = 5} define uma zona de
tolerância: sem ela, pequenas variações no tempo acumulado causariam
oscilações de promoção e rebaixamento. A variável global \texttt{list\_location}, atualizada por
\texttt{pcb\_get\_node\_pid}, comunica o nível atual
do processo para \texttt{pcb\_climb} e \texttt{pcb\_fall}. É uma escolha que evita parâmetros extras, mas cria acoplamento implícito:
uma chamada intercalada a \texttt{pcb\_get\_node\_pid} com outro PID
sobrescreveria \texttt{list\_location} e corromperia a movimentação. Em um
sistema multiprocessado, isso seria uma fonte clássica de \emph{race conditions}.

\subsubsection{Alocação estática e \texttt{time\_elapsed}}

Os nós da lista ligada (que organizam as filas do PCB) são alocados estaticamente em
\texttt{process\_blocks\_nodes[64]}, indexados por PID, impondo o mesmo
limite de 64 processos. O comportamento do escalonador é
funcional, mas vale ressaltar que a função \texttt{time\_elapsed} atualmente simula a 
passagem de tempo, e o feedback de prioridade não reflete ciclos precisos de relógio real.

\subsection{Detalhes da Implementação}

\subsubsection{Estruturas de dados}

\lstset{captionpos=t}
\begin{lstlisting}[language=C, caption={Estruturas do escalonador em \texttt{scheduler.h}}]
struct pcb_entry {
    pll_node *process_list; // cabeca da lista ligada de processos
    unsigned quantum;       // quantum deste nivel
    short process_count;    // numero de processos no nivel
    short next_process;     // indice circular do proximo a ser eleito
};

typedef struct process_node_ll {
    struct process_node_ll *next;
    process proc;
} pll_node;
\end{lstlisting}
\lstset{captionpos=b}

O campo \texttt{next\_process} implementa o revezamento circular dentro de
cada nível: após eleger o processo de índice $i$, o campo avança para
$(i+1) \bmod \texttt{process\_count}$.

\subsubsection{O algoritmo de eleição: \texttt{pcb\_elect}}

Invocada a cada \emph{tick} do timer via \texttt{timer\_callback}, ou sempre que 
um processo cede a CPU voluntariamente (ex: chamando \texttt{waitpid} ou \texttt{exit}),
a função \texttt{pcb\_elect} executa em três fases:

\textbf{Fase 1 — Acumular tempo.} Adiciona o retorno de \texttt{time\_elapsed()}
à variável \path{running_for}.

\textbf{Fase 2 — Avaliar o processo atual.} Se o processo ativo ainda está dentro
do quantum (com a margem) e não mudou de estado, a função retorna sem preemptar:

\lstset{captionpos=t}
\begin{lstlisting}[language=C, caption={Avaliação do quantum em \texttt{pcb\_elect}}]
if (current->state == RUNNING) {
    if (running_for < (quantum + SCHEDULER_QUANTUM_MARGIN))
        return; // ainda dentro do quantum, mantem rodando

    // esgotou o quantum: aplica feedback de prioridade
    if (running_for <= quantum - SCHEDULER_QUANTUM_MARGIN)
        pcb_climb(current_process_node); // usou pouco: promove
    else
        pcb_fall(current_process_node);  // usou muito: rebaixa
        
    running_for = 0;
    current->state = READY;
    current->blocked_by = BT_TIMER;
}
\end{lstlisting}
\lstset{captionpos=b}

\textbf{Fase 3 — Eleger o próximo.} Percorre os 32 níveis do PCB do mais alto para
o mais baixo (nível 0 ao 31). Em cada nível que contenha processos (\texttt{process\_count > 0}),
um laço interno avalia as opções disponíveis garantindo que não entre em loop infinito caso
haja inconsistências:

\lstset{captionpos=t}
\begin{lstlisting}[language=C, caption={Laço de tentativa e eleição de processo}]
for(int tries = 0; tries < pcb[i].process_count; tries++) {
    pll_node *current_node = pll_idx(pcb[i].process_list, pcb[i].next_process);

    if(current_node != NULL && current_node->proc->state == READY) {
        process old_current = current;
        running_for = 0;
        pcb[i].next_process = (pcb[i].next_process + 1) % pcb[i].process_count;

        current = current_node->proc;
        current->state = RUNNING;
        context_switch(&old_current->context, &current->context);
        return;
    }
    pcb[i].next_process = (pcb[i].next_process + 1) % pcb[i].process_count;
}
\end{lstlisting}
\lstset{captionpos=b}

Este laço de tentativas (\texttt{tries}) checa circularmente a fila. 
Ao encontrar um candidato válido em estado \texttt{READY}, atualiza os 
índices de \texttt{current}, altera o estado do novo eleito, 
invoca \texttt{context\_switch} e encerra imediatamente a eleição.

\subsubsection{Promoção e rebaixamento}

\texttt{pcb\_climb} e \texttt{pcb\_fall} movem o processo um nível acima
ou abaixo, respectivamente, removendo o nó do nível atual e inserindo-o no
adjacente. Os limites (nível~0 e nível~31) são respeitados rigidamente: processos nos
extremos não são movidos além deles.


\section{Interface Pública (API)}


\subsection{Módulo de Processos (\texttt{process.h})}

\begin{description}
    \item[\texttt{pid\_t create\_pid()}] Aloca e retorna um PID disponível.
    Retorna \texttt{MAX\_PROCESS\_COUNT} se todos estiverem em uso.

    \item[\texttt{void free\_pid(pid\_t pid)}] Libera o PID para reutilização.

    \item[\texttt{void first\_process(void (*programa)(int, char**))}] Cria e lança
    o processo inicial, passando o ponteiro da função a ser executada como ponto
    de entrada. Deve ser chamada uma única vez pelo \texttt{kmain}. Não retorna.

    \item[\texttt{int fork()}] Cria um processo filho. Retorna o PID do filho
    ao pai e 0 ao filho (wrapper assembly em \texttt{syscall.S}).

    \item[\texttt{void exec(void (*programa)(int, char**), int argc, char** argv)}] Substitui a imagem do
    processo atual pelo novo programa repassando seus argumentos (\texttt{argc} e \texttt{argv}). Não retorna ao chamador.

    \item[\texttt{int waitpid(int pid)}] Bloqueia o processo chamador (\texttt{BLOCKED}) até que o processo filho 
    especificado pelo PID termine, garantindo a coleta do estado Zumbi e liberação da memória. Retorna o PID finalizado.

    \item[\texttt{void exit()}] Encerra o processo atual e avisa ao pai, se este estiver esperando. Não retorna.
\end{description}

\subsection{Módulo do Escalonador (\texttt{scheduler.h})}

\begin{description}
    \item[\texttt{void pcb\_elect(void)}] Seleciona e ativa o próximo
    processo. Chamada automaticamente pelo timer; pode ser invocada
    explicitamente por bloqueios de I/O, \texttt{sys\_waitpid} ou \texttt{sys\_exit}.

    \item[\texttt{void pcb\_insert(unsigned priority\_level, pll\_node* proc)}]
    Insere um processo no nível especificado. Nível~0 é o de maior prioridade.

    \item[\texttt{void pcb\_remove(unsigned priority\_level, pll\_node* proc)}]
    Remove um processo do nível especificado.

    \item[\texttt{pll\_node* pcb\_get\_node\_pid(pid\_t pid)}] Localiza o nó
    do processo pelo PID. Efeito colateral: atualiza \texttt{list\_location}.

    \item[\texttt{void pcb\_climb(pll\_node* proc)}] Move o processo um nível
    acima. Requer chamada prévia de \texttt{pcb\_get\_node\_pid} para o
    mesmo processo.

    \item[\texttt{void pcb\_fall(pll\_node* proc)}] Move o processo um nível
    abaixo. Mesma precondição de \texttt{pcb\_climb}.

    \item[\texttt{pll\_node* pll\_node\_new(process proc)}] Inicializa e
    retorna o nó de lista ligada alocado estaticamente para o processo,
    indexado pelo PID.

    \item[\texttt{int proc\_is\_valid(process proc)}] Retorna verdadeiro se
    o processo não é nulo e não está em estado \texttt{ZOMBIE}.
\end{description}

%-------------------------------------------------------------------------------
\section{Testes e Verificação}
%-------------------------------------------------------------------------------

\subsection{Rotina de Teste}

O \texttt{kmain} inicializa os gerenciadores de memória 
(PMM e VMM) e invoca \path{first_process(prog_init_idle)}. Este processo 
primário cria um filho via \texttt{fork} e executa (\texttt{exec}) a 
\texttt{prog\_shell}, cedendo controle ao usuário através de uma interface interativa.

O teste de validação do módulo de processos e escalonamento consiste em interagir
com essa shell e invocar o comando \texttt{multitask}. Esse programa invoca 
sucessivos \texttt{forks} para criar três processos filhos simultâneos que executam
laços de espera ocupada. Enquanto isso, o processo pai 
utiliza \texttt{waitpid} para aguardar o término sequencial de cada filho. 

\subsection{Saída Esperada}

Ao executar \texttt{make run} e interagir com o sistema, a saída no console 
será análoga à listagem abaixo:

\lstset{captionpos=t}
\begin{lstlisting}[caption={Saída serial do QEMU executando \texttt{multitask}}]
> multitask
Processo pai iniciando...
Filho 1: PID 4
Filho 2: PID 5
Filho 3: PID 6
Pai esperando os filhos...
Esperando 1: PID 4
Filho 1: PID 0
   -> Filho 1 nasceu!
Filho 3: PID 0
   -> Filho 3 nasceu!
Filho 2: PID 0
   -> Filho 2 nasceu!
   <- Filho 1 finalizou!
Esperando 2: PID 5
   <- Filho 3 finalizou!
   <- Filho 2 finalizou!
Esperando 3: PID 6
Todos os filhos terminaram. Pai saindo.
> exit
\end{lstlisting}
\lstset{captionpos=b}

\subsection{Análise da Execução}

A saída comprova o funcionamento integrado de múltiplos subsistemas do kernel:

\begin{enumerate}
    \item \textbf{Criação de Processos (\texttt{fork} e \texttt{exec}):} A shell
    consegue interpretar o comando, realizar o \texttt{fork} e substituir a imagem
    do processo filho pelo código do \texttt{prog\_multitask}. O processo pai, por sua 
    vez, cria três filhos (PIDs 4, 5 e 6).
    
    \item \textbf{Escalonamento e Preempção:} A mensagem \texttt{"Esperando 1: PID 4"} 
    é imediatamente seguida pelas inicializações dos três filhos intercaladas 
    (\texttt{"Filho 1 nasceu!"}, \texttt{"Filho 3 nasceu!"}, \texttt{"Filho 2 nasceu!"}). 
    Isso prova que o timer está interrompendo a CPU ativa e o escalonador MLFQ 
    está realizando as trocas de contexto com sucesso entre os três processos irmãos.
    
    \item \textbf{Sincronização e Coleta (\texttt{waitpid}):} O pai bloqueia em 
    \texttt{waitpid(4)}. Quando o filho 1 finaliza (\texttt{"<- Filho 1 finalizou!"}), 
    o pai é acordado e imediatamente emite \texttt{"Esperando 2: PID 5"}. O processo se repete
    até que todos os filhos terminem e o pai encerre, retornando o controle limpo
    à shell (prompt \texttt{>}), comprovando que a memória e os PIDs foram
    integralmente desalocados sem vazamentos.
\end{enumerate}

% <-- ADICIONE SUA LINHA AQUI

\newpage 

\printbibliography

\end{document}

